\chapter{INTRODUCTION}

\section{Background}
Small and medium-sized trading businesses often operate as intermediaries between suppliers and customers. Their daily operations depend on purchasing goods from suppliers, receiving stock, selling goods to customers, preparing business documents, and monitoring the margin between purchase cost and sales revenue. In such businesses, inventory management is not only a matter of counting product quantities. It also requires the coordination of quotations, purchase orders, sales records, fulfillment status, receivables, payables, and stock availability.

Many small businesses still depend on spreadsheets, paper records, manual document folders, or general-purpose accounting tools. These methods may be workable for a small number of products, but they become difficult to control when the business must track multiple suppliers, product units, transaction statuses, cancelled lines, delayed purchases, customer billing, and supplier payments. Manual workflows also make it difficult to answer operational questions quickly, such as which products are low in stock, which sales can be billed, which purchases are ready for payment, or whether a sale can be fulfilled from available received stock.

The Inventory Management System was developed to address these practical workflow problems for an SME trading environment. The implemented system uses a React frontend, a Django REST Framework backend, and a MySQL relational database. Its design focuses on master data management, purchasing, sales, quotations, billing notes, payment batches, credit notes, dashboard summaries, stock validation, and read-only text queries for selected operational information. The backend is responsible for authoritative validation and stock calculations, while the frontend provides a compact web interface for daily operational work.

The project also includes an AI-assisted inventory chat feature with a limited scope. It is a read-only helper that can summarize supported backend data; it does not create, edit, approve, cancel, or override transactions. Its current supported areas are stock and fulfillment questions, customer or supplier summaries, receivables and payables with exception checks, and reference or line-item lookup.

\section{Objective}
% TODO: The project authors provided this objective section. Review and adjust the final wording if the project scope changes.
\begin{enumerate}
	\item To develop a system that records purchase and sales transactions, including the current status of each process, with an attachment file function that links all supporting documents directly to the corresponding transactions.
	\item To develop an inventory management system that displays the overall product stock in the warehouse and provides tools for calculating essential inventory information.
	\item To integrate an AI chatbot using the ChatGPT API that connects to the system's backend data, allowing users to request supported operational summaries through text queries.
\end{enumerate}

\section{Scope}
The scope of this project is limited to the implemented Inventory Management System codebase. The system covers the following functional areas:

\begin{enumerate}
	\item \textbf{Master data management:} The system manages categories, products, suppliers, and customers. Products support SKUs, previous SKUs, product names, category links, unit conversions, reorder levels, active or disabled status, product pictures, and supplier-specific sourcing information.
	
	\item \textbf{Purchasing workflow:} The system records purchase transactions, supplier information, purchase line items, expected delivery dates, received dates, item statuses, VAT totals, bill discounts, supplier tax invoice references, uploaded documents, and payable totals. Purchase status and purchase item status are synchronized by backend logic.
	
	\item \textbf{Sales workflow:} The system records sales transactions, customer information, customer purchase-order references, line items, delivery statuses, uploaded documents, VAT totals, discounts, billing-note links, credit-note links, and source quotation links. Stock is validated on the server before a sale can move into a stock-deducting status.
	
	\item \textbf{Inventory and stock control:} Available stock is derived from received purchase quantities minus committed sale quantities. Sale allocations use received purchase layers and support first-in, first-out allocation logic. The inventory workspace and dashboard show stock position, reorder information, demand, purchase pipeline, stock value, and product history.
	
	\item \textbf{Quotation workflow:} The system stores quotations with customer and supplier references, quotation validity information, shipping date, payment term, VAT mode, total amounts, normalized line items, and supplier options for quotation items. Quotations can be linked to derived purchases and sales.
	
	\item \textbf{Finance workflow:} The system supports billing notes for customer receivables, payment batches for supplier payables, and credit notes for cancelled or returned sales lines. Eligibility endpoints prevent users from selecting transactions that should not be included in these finance documents.
	
	\item \textbf{Dashboard and reporting:} The backend provides dashboard summary and segment endpoints for operational summaries such as finance, trends, product movement, cashflow, and order coverage.
	
	\item \textbf{Authentication and access flow:} The backend includes JWT login, refresh, and authenticated user profile endpoints. The frontend supports login, token refresh, logout, and guest mode. Backend-wide authentication enforcement is controlled by the \texttt{INVENTORY\_REQUIRE\_AUTH} setting.
	
	\item \textbf{AI-assisted operation:} The system includes a text-based inventory assistant endpoint. Its current scope is read-only support for stock and fulfillment questions, customer or supplier summaries, receivables and payables with exception checks, and reference or line-item lookup.
\end{enumerate}

\section{Expected Results}
The expected result of this project is a functional web-based inventory management system that can record purchase and sales transactions, track the status of each process, and link supporting documents directly to their corresponding transactions. The system is expected to provide an up-to-date view of stock levels based on received purchase quantities and committed sale quantities, together with clear traceability between products, suppliers, customers, purchases, sales, quotations, billing notes, payment batches, and credit notes.

The system is also expected to improve operational accessibility by providing dashboard summaries and a limited AI-assisted text interface. Users should be able to retrieve supported information such as product stock, selected transaction references, receivables, payables, partner summaries, and stock exceptions without manually searching every module. Overall, the completed system is expected to reduce manual checking, improve stock visibility, and support better day-to-day decision making for an SME trading business.

\section{Timeline}
\newcolumntype{L}[1]{>{\raggedright\let\newline\\\arraybackslash\hspace{0pt}}m{#1}}
\newcolumntype{C}[1]{>{\centering\let\newline\\\arraybackslash\hspace{0pt}}m{#1}}
\newcolumntype{R}[1]{>{\raggedleft\let\newline\\\arraybackslash\hspace{0pt}}m{#1}}	

\begin{table}[!ht]
	\footnotesize
	\sloppy
	\centering
	\caption{Project Timeline}
	\label{tab: your-table} %for cross-reference
	\begin{tabular}{|p{2.5cm}|c|c|c|c|c|c|c|c|c|}
		
		\hline
		\multicolumn{1}{|c|}{} & \multicolumn{9}{c|}{\textbf{Timeline}} \\ \cline{2-10}
		\multicolumn{1}{|c|}{} & \multicolumn{4}{c|}{\textbf{2024}} & \multicolumn{5}{c|}{\textbf{2025}} \\ \cline{2-10}
		\multicolumn{1}{|c|}{\multirow{-3}{2cm}{\textbf{Plan}}} & \textbf{Sep} & \textbf{Oct} & \textbf{Nov} & \textbf{Dec} & \textbf{Jan} & \textbf{Feb} & \textbf{Mar} & \textbf{Apr} & \textbf{May} \\ \hline
		1: Initiation & 
		\cellcolor[HTML]{00FF00} & & & & & & & & \\ \hline
		2: Data Analyze & & 
		\cellcolor[HTML]{FFFF00} & & & & & & & \\ \hline
		3: Requirements & & & 
		\cellcolor[HTML]{FFFF00} & & & & & & \\ \hline
		4: Data Setup & & & & 
		\cellcolor[HTML]{FFFF00} & & & & & \\ \hline
		5: Development & & & & & 
		\cellcolor[HTML]{FFFF00} &
		\cellcolor[HTML]{FFFF00} & & & \\ \hline
		6: Integration & & & & & & 
		\cellcolor[HTML]{FFFF00} & & &\\ \hline
		7: User Test& & & & & & & 
		\cellcolor[HTML]{FFFF00} & & \\ \hline
		8: Final Draft & & & & & & & & 
		\cellcolor[HTML]{FFFF00} & \\ \hline
		9:  Submittion & & & & & & & & & 
		\cellcolor[HTML]{FFFF00} \\ \hline
		
	\end{tabular}
\end{table}
